\begin{proof}
(i) By (Z), $\alpha_{4} = g_{5}^{2}/(4\pi L)$, and $g_{5}^{2} = 8\pi^{3}/E_{s}$ gives $\alpha_{4} = 2\pi^{2}/(E_{s}L)$. The depth is E's one length (Premise \claimref{dark-matter.e-equal-mass-per-unit-ball}), counted through once by Premise R3 (\claimref{dark-matter.r3-one-cell-per-unshared-direction}); in units with $\hbar c = 1$, $L = d/\hbar c$. The charged species are those of Proposition \claimref{dark-tower-matter.five-dimensional-content}; their masses do not enter.

(ii) Substituting (i) in (PL), $2\pi^{2}\hbar c/(E_{s}d) \le c\,m_{c}^{2}/\mathrm{TeV}^{2}$, and multiplying by $E_{s}d\,\mathrm{TeV}^{2}/(c\,d_{b})$ gives the inequality. The two values of $K$ are evaluated in rational intervals in the check.

(iii) Setting $m_{c} = 1$ GeV and $d = d_{b}$ leaves $E_{s} \ge K/\mathrm{GeV}^{2}$, and $K/(10^{11}\,\mathrm{GeV}^{3})$ is the factor by which a deposit at that scale exceeds the limit. Setting $E_{s} = 10^{11}$ GeV and $d = d_{b}$ leaves $m_{c}^{2} \ge K/E_{s}$. Along the threshold $m_{c}^{2}E_{s}d$ is constant.

(iv) For a plasma of charge $q$, $\alpha = q^{2}/(4\pi)$, number density $n$ and mass $m_{c}$, $\omega_{p}^{2} = q^{2}n/m_{c} = 4\pi\alpha n/m_{c}$, and $n = \rho/m_{c}$. The growth rate of a streaming instability is a fixed multiple of $\omega_{p}$ at given speeds, so the requirement that it not saturate within a crossing time is $\omega_{p} \le$ a constant, that is $\alpha\rho/m_{c}^{2} \le$ a constant: the limit goes as $m_{c}^{2}/\rho$. E fixes each tower's mass density, whatever its quanta weigh. A five-dimensional tower's share of the dark matter is $(8/3\pi)/((11\pi + 16)/3\pi)$.

(v) $\sigma_{5} \propto \beta^{2}/(\mu v)^{3}$ with $\beta \propto m_{c}g_{5}^{2}$ and $\mu = m_{c}/2$ is proportional to $g_{5}^{4}/(m_{c}v^{3})$, and dividing by $Lm_{c}$ gives $g_{5}^{4}/(m_{c}^{2}v^{3}L)$. The binding condition with $g_{5}^{2} = 8\pi^{3}/E_{s}$ and $C_{F} = 4/3$ is $m_{c} \ge 4\pi^{2}E_{s}/(8\pi^{3}C_{F}) = 3E_{s}/(8\pi)$. For stability: by Proposition \claimref{dark-tower-matter.five-dimensional-content} the centre of the gauge group acts on a state by its triality $k \bmod 3$, which gauge invariance conserves, and the gauge bosons have triality zero. A final state of the same triality as a single coloured particle therefore contains at least one coloured particle, of mass at least that of the lightest coloured species, together with other particles of positive energy; so the lightest coloured species has no decay.

(vi) $(2.8 \times 10^{-14})^{2} \times 10^{6} = 7.84 \times 10^{-22}$, the supplement's coefficient to two figures; dividing by $4\pi$ and repeating (ii) and (iii) gives the stated values.
\end{proof}
